\documentclass{ou}
\begin{document}
\thispagestyle{empty}
\begin{center}
\shortstack{\FGAVO}\vspace{4\baselineskip}
\textbf{РЕФЕРАТ}\\
\setstretch{2}
Информационные системы и~технологии\\
\setstretch{1}
Информационные технологии и информационные системы\vspace{2\baselineskip}\end{center}
\begin{flushleft}
\doublespacing
\begin{tabularx}{\textwidth}{@{} p{0.3\textwidth} l l @{}}
Руководитель & ст. преподаватель & И. Р. Нигматуллин\\
Студенты     &                   &                  \\
\multirow{-2.27}{\hsize}{\singlespacing Понятие и состав информационных технологий как элемента цифровой среды}
             & ГФ25-01Б, 152537671 & А. Е. Корикова\\
             & ГФ25-01Б, 152536528 & Д. А. Васильева\\
             & ГФ25-01Б, 152540016 & И. Н. Морозова\\
\multirow{-2.27}{\hsize}{\singlespacing Структура и функции информационных систем в современном обществе}
             & ГФ25-01Б, 152537684 & О. И. Бочанова\\
             & ГФ25-01Б, 152539986 & М. К. Гинтер\\
             & ГФ25-01Б, 152541471 & А. О. Безруких\\
             & ГФ25-01Б, 152536796 & А. В. Корюкова\\
\multirow{-2.27}{\hsize}{\singlespacing Взаимосвязь между информационными технологиями и информационными системами}
             & ГФ25-01Б, 152540046 & А. Д. Балцевич\\
             & ГФ25-01Б, 152537659 & В. Е. Зайцева\\
             & ГФ25-01Б, 152536592 & А. Ю. Бирюкова\\
\end{tabularx}
\end{flushleft}
\vspace*{\fill}
\begin{center}
\city
\end{center}
\newpage
\thispagestyle{empty}
\setstretch{1}
Продолжение титульного~листа реферата по~теме «Технологии~создания\\и~обработки~графических данных»
\begin{flushleft}
\doublespacing
\begin{tabularx}{\textwidth}{@{} p{0.3\textwidth} l l @{}}
Студенты     &                     &         \\
\multirow{-2.27}{\hsize}{\singlespacing Классификация информационных систем по~сфере применения и уровню автоматизации}
             & ГФ25-01Б, 152540753 & М. И. Бобылев\\
             & ГФ25-01Б, 152539850 & Е. А. Аксаментова\\
             & ГФ25-01Б, 152540765 & А. В. Гайворонская\\
\multirow{-2.27}{\hsize}{\singlespacing Роль информационных систем в управлении данными и поддержке принятия решений}
             & ГФ25-01Б, 152536115 & С. В. Идрисова\\
             & ГФ25-01Б, 152536056 & В. А. Карасова\\
             & ГФ25-01Б, 152539526 & К. Н. Аркадьева\\
             & ГФ25-01Б, 152540131 & Е. Кичко\\
\end{tabularx}
\end{flushleft}
\newpage
\tableofcontents
\newpage
\onehalfspacing
\section{Понятие и состав информационных технологий как элемента цифровой среды}
Характерной чертой нашего времени являются развивающиеся процессы информатизации практически во всех сферах человеческой деятельности. Они привели к формированию новой информационной инфраструктуры, которая связана с~новым типом общественных отношений, с~новой реальностью, с~новыми информационными технологиями различных видов деятельности. 

Современному человеку необходимо знать информационные технологии, уметь успешно применять данные знания при решении как личностных, так и производственных задач повседневной жизни. Информационные технологии меняют нашу жизнь, культуру и профессиональную среду. Они формируют новую информационную среду, где знания становятся доступнее, коммуникации быстрее, а~рабочие процессы -- эффективнее.
\subsection{Понятие информационных технологий}
Информационная технология -- совокупность методов, производственных процессов и технических средств, объединенная технологическим процессом и обеспечивающая сбор, хранение, обработку, вывод и распространение информации для снижения трудоемкости процессов использования ресурсов, повышения их надежности и оперативности. Цель применения ИТ -- производство информации для ее анализа человеком и принятия на~его основе решения по~выполнению какого-либо действия, а~также снижение трудоемкости использования информационных ресурсов.

Методами информационных технологий являются методы сбора, передачи и обработки информации. Совокупность методов и производственных процессов определяет принципы, приемы, мероприятия, регламентирующие проектирование и использование технических средств для обработки данных в предметной области.

Информационные технологии характеризуются следующими основными свойствами:
\begin{itemize}[label=-, itemindent=4em, nosep]
    \item Предметом (объектом) обработки (процесса) являются данные;
    \item Целью процесса является получение информации;
    \item Средствами осуществления процесса являются программные, аппаратные и программно-аппаратные вычислительные комплексы;
    \item Процессы обработки данных разделяются на~операции в соответствии с~данной предметной областью (операция -- это совокупность элементарных действий, выполняемых на~одном рабочем месте, которая приводит к реализации определенной функции обработки данных);
    \item Выбор управляющих воздействий на~процессы должен осуществляться лицами, принимающими решение;
    \item Критериями оптимизации процесса являются своевременность доставки информации пользователю, её надёжность, достоверность, полнота.
\end{itemize}

Информационные технологии играют важную стратегическую роль, так как их применение позволило представить в формализованном виде, пригодном для практического использования, концентрированное выражение научных знаний и практического опыта для реализации и организации социальных процессов. Это привело к экономии затрат труда, времени, энергии, материальных ресурсов, необходимых для осуществления этих процессов.

Информационные технологии в настоящее время получили очень широкое применение. Они используются практически во всех сферах человеческой деятельности. В зависимости от рода предметной области выделяют три класса информационных технологий: глобальные, базовые и~конкретные ИТ.

Глобальные информационные технологии позволяют использовать информационные ресурсы общества в целом:
\begin{enumerate}[itemindent=4em, nosep, label=\arabic*)]
    \item Базовые информационные технологии для определенной области применения (производство, научные исследования, проектирование, обучение и т.д.);
    \item Конкретные информационные технологии реализуют обработку данных при решении конкретных функциональных задач пользователя (например, задачи планирования, учёта, анализа и т.д.);
\end{enumerate}
\subsection{Понятие цифровой среды}
Цифровая среда представляет собой совокупность цифровых технологий, инфраструктуры и данных, которые обеспечивают взаимодействие и обмен информацией в цифровом формате. Это среда, в которой происходит создание, хранение, обработка и передача цифровой информации с~использованием компьютерных и телекоммуникационных систем.

Цифровая среда играет ключевую роль в развитии современного общества и экономики. Она открывает новые возможности для бизнеса, образования, здравоохранения, государственного управления и многих других сфер.

Цифровизация позволяет повысить эффективность, скорость и качество предоставления услуг, а~также расширить доступ к информации и знаниям. Кроме того, цифровая среда способствует развитию инноваций, создает новые рабочие места и стимулирует экономический рост.

В условиях глобализации и технологического прогресса развитие цифровой среды становится ключевым фактором конкурентоспособности и устойчивого развития как отдельных организаций, так и целых стран. Понимание роли и значения цифровой среды в современном мире является важным для успешной адаптации и эффективного использования цифровых технологий.
\subsection{Структупа,принципы реализации и функционирования информационных технологий}

Структура информационных технологий представляет собой устойчивую совокупность взаимосвязанных компонентов, совместная работа которых обеспечивает выполнение цикла обработки информации. Информационные технологии состоят из четырех основных компонентов: 
\begin{enumerate}[itemindent=4em, nosep, label=\arabic*)]
    \item Аппаратное обеспечение (hardware);
    \item Программное обеспечение (software);
    \item Информационные ресурсы;
    \item Организационное обеспечение;
\end{enumerate}

Аппаратное обеспечение (hardware) включает в себя физическое оборудование, необходимое для функционирования информационных технологий.

Устройства обработки данных: центральные процессоры, графические процессоры, серверы, суперкомпьютеры.

Устройства хранения данных: жесткие диски, твердотельные накопители, системы хранения данных, ленточные библиотеки.

Устройства ввода/вывода информации: клавиатуры, манипуляторы, мониторы, принтеры, сканеры, сенсорные панели.

Сетевое и телекоммуникационное оборудование: маршрутизаторы, коммутаторы, модемы, точки доступа.

Компьютерные системы: персональные компьютеры, серверы, а также мобильные устройства.

Периферийные устройства: клавиатуры, мыши, мониторы, принтеры, сканеры, накопители данных.

Программное обеспечение -- это набор инструкций и приложений, которые управляют оборудованием и обеспечивают выполнение задач, определенных пользователем.

Программное обеспечение включает в себя:
\begin{enumerate}[itemindent=4em, nosep, label=\arabic*)]
\item Системное программное обеспечение: операционные системы, драйверы устройств, утилиты, системы виртуализации. Обеспечивает базовое функционирование аппаратной части и управление её ресурсами.
\item Прикладное программное обеспечение: программы, предназначенные для решения задач в конкретных предметных областях, такие как офисные пакеты, системы автоматизированного проектирования, бухгалтерские программы, корпоративные информационные системы.
\item Инструментальное программное обеспечение, или средства разработки: языки программирования, включая Python, Java, C++, системы управления базами данных, средства тестирования, интегрированные среды разработки.
\end{enumerate}

Информационные ресурсы составляют содержательную основу информационных технологий и включает все виды данных и знаний, циркулирующих в системе.

Структурированные данные: реляционные базы данных, электронные таблицы.

Неструктурированные и слабоструктурированные данные: текстовые документы, электронная почта, мультимедийный контент, изображения, аудио и видео.

Базы знаний и онтологии: формализованные модели, используемые в~экспертных системах и системах искусственного интеллекта.

Метаданные: данные, описывающие контекст, содержание и структуру других данных.

Сетевые и телекоммуникационные технологии: обеспечивают единение распределенных компонентов информационных технологий в единую систему.

Каналы связи: проводные, такие как оптоволоконные и медные, и~беспроводные, включая радиоканалы и спутниковые.

Сетевые протоколы: стандартизированные наборы правил для обмена данными, например, стек протоколов TCP/IP, HTTP/HTTPS, FTP, SMTP.

Сетевые технологии и архитектуры: технологии локальных, таких как Ethernet и Wi-Fi, и глобальных сетей, архитектуры «клиент-сервер», одноранговые сети.

Организационное обеспечение включает человеческие ресурсы и регламентирующую документацию, необходимые для эффективного использования и развития информационных технологий.

Персонал: ИТ-специалисты, такие как разработчики, администраторы, аналитики, руководители ИТ-направлений, конечные пользователи.

Организационно-распорядительная документация: регламенты, инструкции, политики информационной безопасности, стандарты оказания ИТ-услуг.

Методологии и системы управления: системы управления ИТ-проектами, управления конфигурациями, инцидентами и изменениями.

\subsection{Принципы реализации информационных технологий}
Основополагающими принципами реализации информационных технологий являются:
\begin{enumerate}[itemindent=4em, nosep, label=\arabic*)]
    \item Системный подход. Рассмотрение ИТ-систем как целостных структур, состоящих из взаимосвязанных элементов. Интеграция аппаратных, программных и информационных ресурсов в единую систему. Обеспечение согласованного взаимодействия компонентов системы;
    \item Модульность. Разделение ИТ-систем на независимые, заменяемые модули. Возможность расширения и модернизации системы путем добавления новых модулей;
    \item Масштабируемость. Возможность наращивания производительности и функциональности ИТ-систем. Обеспечение гибкого распределения ресурсов для покрытия растущих потребностей;
    \item Гибкость и адаптивность. Способность ИТ-систем перестраиваться под изменяющиеся условия и требования. Использование адаптивных алгоритмов и механизмов обратной связи;
    \item Надежность и отказоустойчивость. Обеспечение устойчивой работы ИТ-систем к сбоям, ошибкам и внешним воздействиям. Применение резервирования, репликации, балансировки нагрузки для повышения отказоустойчивости;
    \item Безопасность. Защита ИТ-систем от несанкционированного доступа, утечки данных, вредоносных воздействий. Использование средств аутентификации, шифрования, межсетевых экранов, систем обнаружения вторжений;
\end{enumerate}

\subsection{Функционирование информационных технологий}
Функционирование информационных технологий представляет собой комплекс взаимосвязанных действий и процедур, обеспечивающих сбор, обработку, хранение, передачу и представление информации для решения различных задач в рамках информационных систем. Информационные технологии реализуют следующие основные функции:

\begin{enumerate}[itemindent=4em, nosep, label=\arabic*)]
    \item Сбор и ввод данных. Использование различных источников информации: первичные (наблюдения, эксперименты) и вторичные (документы, базы данных). Преобразование аналоговых данных в цифровую форму с~помощью сенсорных устройств;

    \item Обработка и анализ данных. Применение математических методов обработки информации. Вычислительные технологии: облачные вычисления, высокопроизводительные системы, распределённые вычисления. Реализация аналитических и прогнозных моделей, визуализация данных;

    \item Хранение и накопление данных. Использование различных типов носителей: магнитные, оптические, твердотельные, облачные. Организация баз данных и хранилищ данных для структурированного хранения информации. Обеспечение целостности, доступности и безопасности хранимых данных;

    \item Передача и распространение данных. Реализация каналов передачи данных: локальные, глобальные, мобильные сети;

    \item Представление и визуализация данных. Использование графических, мультимедийных и других форматов для визуального представления информации. Создание интерфейсов и других средств визуализации данных;

    \item Управление и администрирование. Организация информационных процессов и потоков в рамках ИТ-систем. Применение методов и технологий управления жизненным циклом ИТ-систем. Контроль, оптимизация работы информационных технологий;
\end{enumerate}
\subsection{Связь информационных технологий и цифровой среды}
Цифровая среда представляет собой глобальное информационное пространство, формируемое совокупностью, состоящей из электронных средств, телекоммуникационных сетей и информационно-коммуникационных технологий, предназначенное для создания, хранения, обработки, передачи и использования информации в цифровой форме. В рамках этой среды информационные технологии выполняют роль системообразующего элемента, обеспечивающего ее функциональность и целостность.

Роль информационных технологий как элемента цифровой среды раскрывается через выполняемые ими ключевые функции:

1.	Формирование технологической инфраструктуры. Информационные технологии создают фундамент цифровой среды, который включает в себя:

Аппаратно-техническую базу: серверы, центры обработки данных, системы хранения данных, клиентские устройства.

Программно-алгоритмическую базу: операционные системы, системы управления базами данных, платформенное и прикладное программное обеспечение.

Коммуникационную инфраструктуру: сети передачи данных, протоколы обмена, технологии беспроводной связи.

Как отмечается в работах исследователей цифровой экономики, без этой     инфраструктуры существование единой цифровой среды невозможно. Она является базисом для развертывания всех цифровых сервисов и платформ.

2.	Обеспечение информационного взаимодействия. Информационные технологии реализуют процессы сбора, обработки, передачи и представления информации, обеспечивая тем самым коммуникацию между всеми субъектами цифровой среды: физическими лицами, организациями, государственными органами, машинами. Это взаимодействие осуществляется через протоколы и~интерфейсы, что делает среду интегрированной и доступной.

3.	Реализация цифровых сервисов и автоматизация процессов. на~основе информационных технологий создаются и функционируют конкретные сервисы, составляющие пользовательский уровень цифровой среды. К ним относятся:

Государственные информационные системы, например, портал "Госуслуги".

Корпоративные информационные системы, такие как ERP и CRM.

Системы электронного документооборота.

Платформы для дистанционного образования и телемедицины.

Согласно позиции в документах по~цифровой трансформации, именно через эти сервисы цифровая среда оказывает непосредственное влияние на~социально-экономические процессы.

4.	Управление информационными ресурсами. Информационные технологии предоставляют инструменты для работы с~данными как стратегическим активом. Это включает технологии сбора, хранения, анализа, такие как Big Data и Data Mining, и защиты информации. Способность среды эффективно управлять большими объемами данных определяет ее ценность и~потенциал для развития технологий искусственного интеллекта и аналитики.

\section{}
\section{}
\section{Классификация информационных систем по сфере применения и уровню автоматизации}
\subsection{Понятие и сущность информационной системы}
Информационная система -- это совокупность взаимосвязанных методов, программно-технических средств, баз данных и персонала, предназначенная для сбора, передачи, обработки, хранения и выдачи информации. Основная задача информационных систем заключается в обеспечении пользователя актуальными и достоверными данными, необходимыми для принятия решений, анализа ситуации и управления различными процессами.

Структура любой информационной системы включает несколько ключевых компонентов. Во-первых, это аппаратное обеспечение, которое представляет собой вычислительную технику: персональные компьютеры, серверы, сетевое оборудование, устройства ввода и вывода информации. Во-вторых, программное обеспечение, обеспечивающее работу системы и её функциональные возможности: операционные системы, прикладные программы, сервисы и инструменты обработки данных.

Дополняют систему организационное обеспечение (инструкции, регламенты, правила работы), информационное обеспечение (базы данных, классификаторы, справочники, наборы знаний) и персонал, то есть специалисты, которые проектируют, обслуживают и используют систему.

Особенность любой информационной системы состоит в том, что она создаётся под конкретные задачи организации или сферы деятельности. Именно поэтому классификация по сфере применения позволяет выделить группы систем с учётом их функций.
\subsection{Классификация информационных систем по сфере применения}
По сфере применения информационные системы делятся на несколько крупных категорий, каждая из которых ориентирована на решение задач.

Управленческие ИС используются для поддержки процессов управления на уровне организаций и предприятий. Они собирают и анализируют данные о~деятельности, формируют отчёты, помогают в планировании и контроле.

К числу распространённых систем относятся:
\begin{itemize}[label=-, itemindent=4em, nosep]
\item MIS (Management Information Systems) -- системы информационной поддержки менеджмента, обеспечивающие доступ к данным для среднего управленческого звена;
\item EIS (Executive Information Systems) -- системы для руководителей высшего звена, предоставляющие стратегическую информацию;
\item ERP (Enterprise Resource Planning) -- системы планирования ресурсов предприятия, интегрирующие деятельность всех подразделений.
\end{itemize}

Использование подобных систем обеспечивает единое информационное пространство и повышает эффективность управленческих решений.

Производственные информационные системы нужны для~контроля производственных процессов, обеспечения качества продукции и повышения эффективности производства.

К основным их типам относятся:
\begin{itemize}[label=-, itemindent=4em, nosep]
\item SCADA-системы, обеспечивающие мониторинг и управление технологическими процессами;
\item MES-системы (Manufacturing Execution Systems), которые управляют производственными операциями в режиме реального времени;
\item APS-системы (Advanced Planning and Scheduling), выполняющие планирование и оптимизацию производственных задач.
\end{itemize}

Производственные ИС позволяют существенно снизить количество ошибок, оптимизировать ресурсы и повысить производственную дисциплину.

Финансово-экономические информационные системы применяются в~банках, страховых компаниях, бухгалтерии и инвестиционной сфере.

К ним относятся:
\begin{itemize}[label=-, itemindent=4em, nosep]
\item Автоматизированные банковские системы;
\item Программы бухгалтерского учёта (например, 1С);
\item Аналитические инвестиционные платформы.
\end{itemize}

Основная задача данных систем -- точность финансовых операций, снижение рисков и автоматизация расчётных процессов.

В торговле информационные системы применяются для учёта товаров, анализа спроса, управления складом и автоматизации кассовых операций.

Среди наиболее распространённых:
\begin{itemize}[label=-, itemindent=4em, nosep]
\item POS-системы, обеспечивающие работу касс;
\item CRM-системы, предназначенные для взаимодействия с клиентами;
\item Системы складского и логистического учёта.
\end{itemize}

Они обеспечивают оперативность обмена данными, улучшение обслуживания и прозрачность товарных потоков.

Медицинские ИС помогают вести электронные медицинские карты, хранить большие объёмы медицинских данных, организовывать работу врачей и поддерживать диагностические решения.

К ним относятся:
\begin{itemize}[label=-, itemindent=4em, nosep]
\item Системы электронных медицинских карт;
\item Телемедицинские системы;
\item Системы поддержки диагностических решений (CDSS).
\end{itemize}

Образовательные ИС применяются для автоматизации учебного процесса, управления учебным заведением и проведения дистанционного обучения.

В их числе:
\begin{itemize}[label=-, itemindent=4em, nosep]
\item LMS (Learning Management Systems) -- системы дистанционного обучения;
\item Электронные журналы и дневники;
\item Системы тестирования и контроля знаний.
\end{itemize}

Такие решения делают образование более доступным и повышают эффективность педагогической деятельности.

Государственные ИС обеспечивают автоматизацию взаимодействия между государственными органами, гражданами и организациями.

Примеры:
\begin{itemize}[label=-, itemindent=4em, nosep]
\item Порталы государственных услуг;
\item Системы электронного документооборота;
\item Государственные реестры и кадастры.
\end{itemize}

Они увеличивают прозрачность работы государственных органов и~сокращают время предоставления услуг.

Для научной деятельности создаются сложные информационные системы, позволяющие анализировать большие массивы данных, моделировать процессы и поддерживать публикационную активность.

К ним относятся:
\begin{itemize}[label=-, itemindent=4em, nosep]
\item Системы научных публикаций;
\item Платформы обработки больших данных (Big Data);
\item Системы моделирования (например, MATLAB, ANSYS).
\end{itemize}
\subsection{Классификация информационных систем по уровню автоматизации}
Ручные системы основаны на обработке информации человеком без~использования компьютерных технологий. В настоящее время они встречаются редко, но могут применяться при отсутствии технических средств.

Для них характерны низкая скорость обработки данных, высокая вероятность ошибки и недостаточная надёжность.

В автоматизированных системах значительная часть операций выполняется техническими средствами, однако человек продолжает принимать основные решения и контролировать процесс.

Такие системы позволяют автоматизировать рутинные операции, но~требуют подготовки персонала и постоянного участия специалистов.

Автоматические ИС функционируют без участия человека: они самостоятельно собирают данные, обрабатывают их, принимают решения и формируют результаты. Участие человека требуется только в аварийных ситуациях или для общего контроля.

К ним относятся системы автоматического регулирования, охранные системы, автономные комплексы мониторинга.

Интеллектуальные ИС основаны на методах искусственного интеллекта и анализе больших объёмов данных. Они способны самостоятельно выявлять закономерности, прогнозировать события, адаптироваться к изменениям.

Примеры: рекомендательные сервисы, экспертные системы, системы прогнозирования.
\subsection{Сравнительный анализ классификаций}
Анализ классификаций показывает, что выбор вида информационной системы зависит от характера задач и требований к автоматизации.

Например, в здравоохранении использование полностью автоматических систем недопустимо при диагностике, поэтому здесь преобладают решения, где человек остаётся ключевым звеном. В производстве многие процессы могут выполнять автоматические системы, так как они требуют высокой точности и скорости реакции. В государственном управлении важно сохранять управленческий контроль, поэтому здесь широко применяются автоматизированные ИС.

В научной сфере, где важно анализировать большие массивы данных, активно используются интеллектуальные системы, способные выполнять сложную обработку информации. Таким образом, степень автоматизации напрямую связана с требованиями к надёжности, рискам и характером деятельности.
\section{}
\newpage
{
\sloppy
\nocite{*}
\bibliographystyle{plainurl}
\bibliography{report1}
}
\end{document}